\section{Introduction}
\label{sec:intro}

Shine a laser on a mineral and a tiny fraction of the scattered light comes back shifted in color by amounts that depend on how the mineral's atoms vibrate. That pattern of Raman shifts works like a fingerprint for the crystal structure. Raman spectroscopy has therefore become a standard way to identify minerals, and laboratory libraries such as RRUFF \citep{lafuente2015rruff} let anyone match an unknown spectrum against thousands of reference minerals. Since 2021, NASA's Perseverance rover has carried two Raman instruments to Jezero crater \citep{farley2020mars2020}. SuperCam measures rocks several meters away with a pulsed green laser \citep{maurice2021supercam,wiens2021supercam}, and SHERLOC uses a deep-ultraviolet laser at close range \citep{bhartia2021sherloc}. SuperCam's Raman spectra have helped identify olivine, carbonates, sulfates, and perchlorates on the crater floor and beyond (e.g., \citep{farley2022aqueously,liu2022olivine,wiens2022stratified,beyssac2023seitah,clave2023carbonate,lopezreyes2025supercam}).

At present, these identifications are made by expert scientists who look at each spectrum \citep{jpl2026supercamsis}. Automating that step with machine learning is attractive for the future, both to keep up with growing data volumes and to give rovers more autonomy. A long line of work trains classifiers on laboratory spectra and reports high accuracies, using library matching \citep{carey2015machine,berrada2025ramanmatch}, classical machine learning \citep{ishikawa2013automated,jahoda2021machine}, and convolutional neural networks (CNNs) \citep{liu2017deep,soysal2025reevaluating}. Benchmarks keep finding, however, that models trained on one dataset or instrument rarely generalize to another \citep{koddenbrock2026ramanbench,sineesh2026benchmarking,qin2026perturbation}. Mars-oriented studies have so far stayed mostly in the laboratory, using lab instruments built to mimic SuperCam \citep{manrique2024simulcam,veneranda2023chemometrics,veneranda2023carbonate,julvegonzalez2023machine,johnsen2025dualband}. \citet{berlanga2022cnn} trained a CNN on lab spectra, including dust-covered rocks, and ran a preliminary, qualitative check on a few early SuperCam spectra. \citet{wang2026sherloc} trained and tested deep models entirely on expert-labeled SHERLOC flight spectra \citep{corpolongo2023sherloc}, and \citet{clave2024unsupervised} decomposed the SuperCam Raman dataset without any lab training. For laser-induced breakdown spectroscopy (LIBS), the other main laser technique on Mars rovers, moving models from the lab to Mars is an established research topic \citep{cui2025libs}. For Raman, we found no quantitative test of how a lab-trained classifier performs on labeled SuperCam spectra. That gap matters: a classifier that is good on held-out \emph{laboratory} spectra is not necessarily good on \emph{flight} spectra. Real rover data differ in resolution, noise, background, and instrument response, and they usually record mixtures of minerals rather than pure samples.

We focus on SuperCam because its 532\,nm laser matches the excitation used for much of RRUFF, whereas SHERLOC's deep-ultraviolet spectra need their own reference libraries \citep{razzellhollis2021library}.

This paper asks a simple question: \textbf{if you build a ``Shazam for minerals'' from Earth lab spectra, does it work on Mars?} We answer it using public data only:
\begin{itemize}
\item We build a reproducible benchmark that links the SuperCam team's published mineral identifications to the matching calibrated SuperCam Raman spectra in the Planetary Data System \citep{wiens2021supercambundle}. It contains \nMarsMatched{} spectra, and \nMarsEval{} of them from \nMarsTargets{} rocks are held out for testing.
\item We show that standard library search, which is excellent in the lab (\labGrpCleanCosineOpen{} mineral-group accuracy on held-out lab samples), fails on Mars when any mineral is allowed. It reaches \balCosineOpen{} balanced accuracy, and \fracAmorphOpen{} of Mars spectra are matched to glasses or other amorphous materials.
\item We trace this failure mainly to the instrument's broad ``fiber bump'' below about 560\,cm$^{-1}$ \citep{clave2024unsupervised,lopezreyes2025supercam}. When predictions are restricted to minerals known from Jezero and that region is simply skipped, library search reaches \balCosineWindowMars{} balanced accuracy, the highest score in this study.
\item We find a clear \emph{sim-to-real gap}: training a CNN on randomized ``simulated Mars'' versions of lab spectra makes it the most robust model on degraded lab spectra, yet \emph{lowers} its accuracy on real Mars spectra (\balCnnMarsifiedMars{} vs.\ \balCnnCleanMars{} without the simulation). Tested one at a time, every simulated effect hurts, and synthetic mineral mixing hurts most.
\end{itemize}
The accompanying code regenerates every number in this paper from the public archives.
